\documentclass[11pt]{article}

\usepackage{amsmath}
\usepackage{amssymb}
\usepackage{array}
\usepackage{booktabs}
\usepackage{bm}
\usepackage{geometry}
\usepackage{hyperref}
\usepackage{xcolor}

\geometry{a4paper, margin=1in}
\hypersetup{colorlinks=true, linkcolor=blue!50!black, citecolor=blue!50!black, urlcolor=blue!50!black}

\newcommand{\stream}{STREAM}
\newcommand{\uce}{UCE}
\newcommand{\cre}{cCRE}
\newcommand{\E}{\mathbb{E}}
\newcommand{\R}{\mathbb{R}}
\newcommand{\N}{\mathcal{N}}
\newcommand{\dd}{\mathrm{d}}

\title{Proposal: stochastic population dynamics for STREAM}
\author{}
\date{}

\begin{document}

\maketitle

\begin{abstract}
The current \stream{} model learns a deterministic, autonomous expression-space
vector field from optimal-transport-paired developmental snapshots. Across
balanced, partial, and unbalanced transport; PCA, KNN, and metacell denoising;
adjacent and skipped intervals; and contiguous held-out blocks, causal endpoint
prediction remains close to expression persistence. This proposal replaces the
deterministic ordinary differential equation with a neural stochastic
differential equation and trains predicted populations directly against
observed endpoint populations. The drift remains a per-gene,
regulatory-sequence-conditioned velocity. A score head and prescribed
diffusion represent branching and unresolved fate variation. The recommended
first objective is simulation-free score-and-flow matching on noisy stochastic
interpolants sampled from an entropic endpoint coupling. A complementary
population-level objective directly minimizes endpoint Sinkhorn divergence by
reparameterized SDE rollout. We define both objectives, their assumptions,
regularization, evaluation, implementation stages, and decision criteria for a
focused stochastic-flow ablation.
\end{abstract}

\section{Motivation}

The deterministic autonomous model has the form
\begin{equation}
  \frac{\dd x}{\dd t}
  = v_\theta\!\left(E(x), \{Z_g\}_{g=1}^{G}\right),
  \label{eq:ode}
\end{equation}
where $x\in\R^G_{\geq 0}$ is expression, $E$ is the frozen \uce{}
encoder, and $Z_g$ contains the AlphaGenome-derived regulatory tokens linked to
gene $g$. The model predicts one expression velocity per gene. Training uses
conditional flow-matching targets formed from optimal-transport pairs between
cross-sectional timepoints.

The completed ablations do not establish meaningful improvement over
persistence. On the original two-stage holdout, the best mean
persistence-normalized Sinkhorn skill is approximately $0.009$, while
gene-level mean-shift $R^2$ remains negative. Holding out the middle three
stages together increases the best mean Sinkhorn skill to approximately $0.02$,
but mean-shift $R^2$ is still negative. A long observed-endpoint chord across
the held-out block improves some mean-shift errors without producing a
consistent distributional improvement.

An initial score-and-flow implementation with an autonomous flow head and an
independent $\tau$-conditioned score head also failed. Validation denoising MSE
remained approximately one, the null value for standardized Gaussian noise,
under raw, KNN-smoothed, and metacell-smoothed targets. Diffusion improved
Sinkhorn distance while worsening gene-level mean shifts, but this could not be
attributed to a learned score. This result motivates the coupled
conditional-moment parameterization below and an explicit zero-score diffusion
control.

These results suggest that the main limitation is not simply model capacity.
Cross-sectional OT does not identify a unique descendant for each source cell,
development branches, and cell numbers change through proliferation and death.
A deterministic map can only assign one local velocity to a state, whereas the
data may support a distribution of descendants. Moreover, pairwise velocity
regression is not the same objective as accurate endpoint population
prediction.

\section{Objective and scope}

The primary question is whether representing conditional fate variability and
training against endpoint distributions improves causal forecasting over
persistence. The first stochastic model should therefore:
\begin{itemize}
  \item retain gene-valued drift and the existing \uce{}/\cre{} conditioning;
  \item exclude absolute developmental-time labels from the state model;
  \item compare a strictly autonomous drift with a bridge model conditioned
        only on the local interpolation coordinate $\tau$;
  \item avoid requiring one-to-one cross-time cell pairs in the primary loss;
  \item generate multiple descendants per source cell;
  \item optimize the same type of endpoint discrepancy used for evaluation;
  \item separate stochastic biological variation from unrestricted technical
        noise through a constrained diffusion parameterization.
\end{itemize}

The initial experiment will not attempt to model division events explicitly,
infer lineage trees, or reproduce raw count sampling. A growth-weighted
extension is described as a second stage.

\section{Autonomous neural SDE}

We propose the expression-space stochastic differential equation
\begin{equation}
  \dd X_t
  = \mu_\theta(X_t)\,\dd t
  + G_\phi(X_t)\,\dd W_t,
  \qquad X_t\in\R^G_{\geq 0},
  \label{eq:sde}
\end{equation}
where $W_t$ is an $r$-dimensional Wiener process. The drift is the existing
regulatory model,
\begin{equation}
  \mu_{\theta,g}(x)
  = \operatorname{STREAM}_{\theta,g}
    \!\left(E(x), Z_g\right).
  \label{eq:drift}
\end{equation}
Thus stochastic training does not change the interpretation of the predicted
gene velocity.

\subsection{Low-rank diffusion}

A dense $G\times G$ diffusion matrix is neither identifiable nor tractable for
$G=10{,}000$. Fully independent gene noise is computationally simple but fails
to represent coordinated lineage changes. We instead propose a low-rank
diffusion with $r=8$--$32$:
\begin{equation}
  G_\phi(x)
  =
  \begin{bmatrix}
    b_{\phi,1}(E(x),Z_1)^\top \\
    \vdots \\
    b_{\phi,G}(E(x),Z_G)^\top
  \end{bmatrix}
  \in \R^{G\times r}.
  \label{eq:lowrank}
\end{equation}
Each gene's promoter-token representation receives a shared $r$-dimensional
diffusion head. This preserves the shared regulatory architecture and allows
one latent noise direction to alter a coordinated set of genes. The diffusion
head can use a bounded output,
\begin{equation}
  b_{\phi,g}=b_{\max}\tanh(\widetilde b_{\phi,g}),
\end{equation}
which prevents the model from improving a distributional loss by adding
unbounded noise.

The lowest-risk first implementation may instead use a state-conditioned
scalar $a_\phi(E(x))$ multiplying a learned global $G\times r$ loading matrix.
This is less expressive but provides a useful test of whether any coordinated
stochasticity helps before adding a second regulatory output head.

\section{Differentiable stochastic rollout}

For an interval of length $T=t_1-t_0$, Euler--Maruyama with $K$ steps gives
\begin{align}
  X^{(m)}_{k+1}
  &= \Pi_+\!\left[
       X^{(m)}_k
       + \mu_\theta(X^{(m)}_k)\Delta t
       + G_\phi(X^{(m)}_k)\sqrt{\Delta t}\,
         \epsilon^{(m)}_k
     \right], \\
  \epsilon^{(m)}_k &\sim \N(0,I_r),
  \qquad \Delta t=T/K,
  \label{eq:em}
\end{align}
where $m$ indexes stochastic particles and $\Pi_+$ clamps negative expression
values to zero. The Gaussian draws are external random inputs, so
Equation~\ref{eq:em} is reparameterized and gradients propagate to drift and
diffusion parameters.

Online \uce{} preprocessing includes discrete gene sampling and is frozen.
Consequently, the current encoder should be treated as a stop-gradient state
map at each step. Gradients still reach earlier updates through the additive
identity path in Equation~\ref{eq:em}, and each step contributes direct
parameter gradients. This avoids claiming differentiability through discrete
\uce{} tokenization. Fixed random seeds should be used within each optimization
step to reduce gradient variance.

\section{Recommended first approach: score-and-flow matching}

Simulation-free score-and-flow matching avoids backpropagation through
multi-step stochastic rollouts and repeated online \uce{} evaluations. Given
coupled endpoints $(x_0,x_1)$, define the noisy stochastic interpolant
\begin{equation}
  X_\tau
  = (1-\tau)x_0+\tau x_1+\gamma(\tau)\Sigma z,
  \qquad z\sim\N(0,I),
  \qquad \tau\sim\operatorname{Uniform}(0,1),
  \label{eq:stochastic-interpolant}
\end{equation}
where $\gamma(0)=\gamma(1)=0$. One possible schedule is
$\gamma(\tau)=\sigma\sqrt{\tau(1-\tau)}$, with endpoint clipping or a smooth
bounded alternative used to avoid numerical singularities.

Here $\Sigma$ is a fixed diagonal matrix of training-only per-gene expression
scales. It makes perturbations comparable across genes while retaining
gene-valued states and outputs.

\subsection{Coupled conditional-moment heads}

The flow and score must not be unrelated readouts. Define two shared
$\tau$-conditioned conditional moments,
\begin{align}
  u_\theta(x,\tau)
  &\approx \E[(x_1-x_0)/\Delta t\mid X_\tau=x], \\
  e_\theta(x,\tau)
  &\approx \E[z\mid X_\tau=x].
  \label{eq:conditional-moments}
\end{align}
Both are predicted from the same time-modulated promoter features. The same
noise estimate $e_\theta$ then constructs both fields:
\begin{align}
  b_\theta(x,\tau)
  &=u_\theta(x,\tau)
    +\frac{\dot\gamma(\tau)}{\Delta t}\Sigma e_\theta(x,\tau), \\
  s_\theta(x,\tau)
  &=-\Sigma^{-1}\frac{e_\theta(x,\tau)}{\gamma(\tau)}.
  \label{eq:coupled-fields}
\end{align}
Thus the noise contribution to the flow and the score cannot disagree without
increasing the same noise-regression loss. The objective is
\begin{equation}
  \mathcal L_{\mathrm{coupled}}
  =\E\left[
    \left\|\frac{b_\theta(X_\tau,\tau)-\dot X_\tau}{v_{\mathrm{rms}}}\right\|_2^2
    +\lambda_s\left\|e_\theta(X_\tau,\tau)-z\right\|_2^2
  \right],
  \label{eq:coupled-loss}
\end{equation}
where $\dot X_\tau=(x_1-x_0)/\Delta t+\dot\gamma(\tau)\Sigma
z/\Delta t$. Predicting standardized noise is equivalent to a
$\gamma(\tau)^2$-weighted score loss and avoids endpoint singularities.
The normalizer $v_{\mathrm{rms}}$ combines the training-only endpoint-velocity
RMS with the analytic RMS of the bridge-noise term, preventing intermediate
noise levels from dominating the coupled loss.

The same checkpoint contains a strictly autonomous control,
\begin{equation}
  b^{\mathrm{aut}}_{\theta,g}(x)
  =\operatorname{STREAM}^{\mathrm{aut}}_{\theta,g}(E(x),Z_g),
\end{equation}
trained against $(x_1-x_0)/\Delta t$ before applying the $\tau$ modulation.
This isolates the value of the local bridge coordinate from architecture and
endpoint samples. Absolute developmental time is supplied to neither path.

\subsection{SDE induced by the matched fields}

If $b_\theta$ satisfies the continuity equation and $s_\theta$ estimates the
marginal score, then the fields define the SDE
\begin{equation}
  \dd X_t
  = \left[
      b_\theta(X_t,\tau(t))
      + \kappa(t)\Sigma^2s_\theta(X_t,\tau(t))
    \right]\dd t
    + \sqrt{2\kappa(t)}\Sigma\,\dd W_t,
  \qquad \tau(t)=\frac{t-t_0}{t_1-t_0}.
  \label{eq:score-flow-sde}
\end{equation}
The score correction offsets diffusion in the Fokker--Planck equation, so an
exact model has the prescribed interpolant marginals for any nonnegative
diffusion schedule $\kappa$. Varying $\kappa$ therefore controls stochasticity
without changing the ideal endpoint distribution. Setting $\kappa=0$ recovers
the probability-flow ODE and provides a matched deterministic control.
For every positive $\kappa$, evaluation also sets the score correction to zero
while retaining identical Brownian draws. This distinguishes a learned-score
benefit from generic variance broadening.

\subsection{Entropic endpoint coupling}

Equation~\ref{eq:stochastic-interpolant} still requires a joint endpoint
distribution. Independent pairing creates unnecessarily broad paths, while a
nearly deterministic OT plan recreates the current arbitrary-pair problem. The
recommended compromise is an entropic OT or KL-relaxed UOT coupling. Sampling
many endpoint pairs from its nonzero support exposes the model to multiple
plausible descendants. This is the simulation-free Schrodinger bridge matching
construction described in related work
\cite{albergo2023stochastic,tong2024simulationfree}.

This approach is simulation-free during fitting: sample an endpoint pair,
$\tau$, and $z$, then regress analytic flow and score targets. It is therefore
well suited to frozen, partly discrete online \uce{} preprocessing. Its main
cost is conceptual: the network must receive $\tau$, because bridge noise and
drift change along the interpolation. Here $\tau$ is a local mathematical path
coordinate, not an absolute developmental stage label. Nevertheless, a model
conditioned on $\tau$ is not strictly autonomous in the sense used by the
current no-time-conditioning experiments.

\section{Complementary population-level objective}

Let $\{x_i^0\}_{i=1}^{B}$ be source cells and
$\{x_j^1\}_{j=1}^{B_1}$ an independently sampled target population. Generate
$M$ particles from every source cell and collect the predicted endpoint set
\begin{equation}
  \widehat{\mathcal X}_1
  = \left\{X_T^{(i,m)}: i=1,\ldots,B,\;m=1,\ldots,M\right\}.
\end{equation}
The principal loss is a debiased Sinkhorn divergence in a train-only,
log-normalized PCA scoring space $P$:
\begin{equation}
  \mathcal L_{\mathrm{pop}}
  = S_\varepsilon\!\left(
      P(\widehat{\mathcal X}_1),
      P(\mathcal X_1)
    \right).
  \label{eq:sinkhorn}
\end{equation}
PCA is used here as a differentiable population metric and denoising transform,
not as the model's dynamical output space. Drift and simulated states remain
gene-valued.

This loss aligns training with held-out endpoint evaluation and removes the
need to select an individual OT partner for every cell. A minibatch size of
$B=128$--$256$, $M=4$ particles, $K=4$ rollout steps, and rank $r=8$ is an
appropriate initial configuration. Larger $M$, $K$, and $r$ should be tested
only after establishing a positive endpoint signal.

\subsection{Regularization}

Without constraints, diffusion can absorb model error or produce an overly
broad population. We propose
\begin{equation}
  \mathcal L
  = \mathcal L_{\mathrm{pop}}
  + \lambda_G\E\left[\lVert G_\phi(X)\rVert_F^2\right]
  + \lambda_\mu\E\left[\lVert \mu_\theta(X)\rVert_2^2\right]
  + \lambda_{\mathrm{mom}}\mathcal L_{\mathrm{mom}},
  \label{eq:full-loss}
\end{equation}
where $\mathcal L_{\mathrm{mom}}$ penalizes large errors in endpoint gene means
and variances. The moment term is important because the completed deterministic
experiments occasionally improve Sinkhorn distance while predicting incorrect
gene-level mean shifts. Diffusion should be initialized near zero, and
$b_{\max}$ should be calibrated against within-timepoint variation.

To keep the distributional loss informative, target and source samples should
be drawn independently at every optimization step. Validation must use fixed
cell splits, seeds, source cells, target cells, and Brownian increments.

\section{Cheaper paired-transition baseline}

Before full rollout training, the existing OT pair banks can support a
heteroscedastic transition likelihood:
\begin{equation}
  x_1-x_0
  \sim
  \N\!\left(
    \mu_\theta(x_0)\Delta t,
    \Sigma_\phi(x_0)\Delta t
  \right).
  \label{eq:nll}
\end{equation}
For diagonal $\Sigma_\phi$, the negative log likelihood is inexpensive and
tests whether state-dependent uncertainty improves calibration. However, this
objective inherits arbitrary OT pairing and can reward inflated variance. It
should be treated as an implementation smoke test and baseline, not the main
stochastic-flow result.

\section{Other alternatives}

\paragraph{Iterative Schrodinger bridge fitting.}
Iterative proportional fitting can learn forward and backward stochastic
bridges between endpoint distributions~\cite{debortoli2021diffusion}. This is a
principled distributional approach but requires repeated simulation and
time-dependent score or drift models. It should be considered if the simpler
autonomous SDE shows a signal.

\paragraph{Independent per-gene noise.}
A diagonal diffusion head is easy to implement but mostly represents technical
noise and cannot express correlated fate decisions. It remains a useful
negative control against low-rank diffusion.

\section{Growth and changing population mass}

An SDE transports probability mass but does not create or remove it. If
proliferation and death are central, add a growth-rate head
$\gamma_\psi(E(x))$ and particle weights
\begin{equation}
  w_{k+1}=w_k\exp\!\left(\gamma_\psi(X_k)\Delta t\right).
\end{equation}
Training then uses weighted or unbalanced Sinkhorn divergence. This separates
state motion, stochastic fate dispersion, and mass change. Growth should be
introduced only after comparing deterministic and stochastic transport under
fixed mass; otherwise diffusion and growth may compensate for each other and
be difficult to identify.

\section{Evaluation}

For each source cell, generate $M_{\mathrm{eval}}=16$--$64$ particles with
fixed seeds. Report:
\begin{itemize}
  \item debiased Sinkhorn divergence and skill relative to persistence;
  \item gene-level mean-shift $R^2$ and mean absolute error;
  \item mean per-gene Wasserstein-1 distance;
  \item endpoint covariance error in train-only PCA space;
  \item coverage and calibration of observed cells under generated particles;
  \item metrics stratified by annotated cell type or lineage where available;
  \item the autonomous ODE and coupled probability-flow ODE from the same
        checkpoint;
  \item matched stochastic rollouts with learned versus zero score correction
        and identical Brownian seeds;
  \item sensitivity to particle count and integration steps.
\end{itemize}

Persistence should be represented by repeated source cells, so it has the same
number of particles as the stochastic model. A deterministic \stream{} model,
the same SDE with $G=0$, and a pure-diffusion model with $\mu=0$ provide the
necessary ablations. Pairwise displacement $R^2$ should remain secondary
because no individual descendant is observed.

\section{Experimental plan}

\begin{table}[ht]
\centering
\caption{Proposed staged experiments. Advancement requires improved held-out
population metrics without loss of gene-level mean-shift accuracy.}
\small
\setlength{\tabcolsep}{4pt}
\begin{tabular}{p{0.06\linewidth}p{0.22\linewidth}p{0.28\linewidth}p{0.28\linewidth}}
\toprule
Stage & Model & Training objective & Decision \\
\midrule
0 & Existing deterministic drift & Existing CFM loss & Reproduce current checkpoint and evaluation \\
1 & Score-and-flow STREAM & Entropic-coupling flow and denoising-score regression & Verify analytic targets and endpoint marginals \\
2 & Matched SDE and ODE & Vary $\kappa$; disable score at inference & Require positive skill and a specific stochastic benefit \\
3 & Global rank-8 autonomous diffusion & Population Sinkhorn, $M=4$, $K=4$ & Compare simulation-free and rollout-trained objectives \\
4 & CRE-conditioned rank-8 diffusion & Population Sinkhorn plus moment loss & Test state/gene-specific uncertainty \\
5 & Optional growth head & Weighted unbalanced Sinkhorn & Proceed only if mass mismatch remains limiting \\
\bottomrule
\end{tabular}
\end{table}

The first full comparison should use the contiguous middle-three-stage holdout.
Train separate adjacent, skip-1, and skip-2 models only if Stage 1 beats
persistence. The observed bridge interval across the block can be added as a
subsequent controlled factor. Because Stage 1 is time-dependent while Stage 3
is autonomous, conclusions must distinguish improved stochastic modeling from
the additional interpolation-coordinate input.

\section{Implementation plan}

The implementation should preserve package boundaries:
\begin{enumerate}
  \item Add an autonomous velocity readout and shared $\tau$-conditioned
        velocity/noise readouts to the reusable model package. Construct flow
        and score analytically from the same noise prediction. Keep current
        drift-only checkpoints backward compatible.
  \item Add a stochastic-interpolant sampler returning entropically coupled
        endpoints, $\tau$, Gaussian noise, and analytic flow/score targets.
  \item Add the coupled flow/noise loss, with tests against analytic Gaussian
        bridges, known marginal scores, and the zero-score SDE control.
  \item Add a reusable SDE rollout for Equation~\ref{eq:score-flow-sde}, plus a
        probability-flow ODE mode with $\kappa=0$.
  \item Add an autonomous drift--diffusion output path and a reusable
        Euler--Maruyama rollout accepting fixed Gaussian noise,
        particle count, diffusion rank, and nonnegativity projection.
  \item Add differentiable train-only PCA scoring and debiased Sinkhorn loss in
        torch. Do not route model outputs through PCA.
  \item Add an interval-stratified population sampler that returns independent
        source and target cells rather than OT pairs.
  \item Add fixed stochastic validation batches containing cell indices and
        Brownian seeds. Select checkpoints by endpoint Sinkhorn EMA, with
        moment metrics as safeguards.
  \item Extend causal evaluation to generate multiple particles, report
        stochastic calibration, and retain deterministic/persistence baselines.
  \item Run synthetic tests where a branching SDE has known drift and diffusion
        before fitting atlas data.
\end{enumerate}

The highest-risk engineering component is repeated online \uce{} evaluation
for $B\times M$ particles. Initial rollout training should cache repeated source
encodings where possible, use small $M$ and $K$, and profile whether \uce{} or
the gene-conditioned drift dominates runtime. A lower-cost expression-state
SDE can validate the training objective before enabling online \uce{}.

\section{Success and stopping criteria}

The stochastic model should be considered useful only if it satisfies all of
the following on held-out stages:
\begin{itemize}
  \item Sinkhorn skill is consistently positive across horizons and materially
        larger than the current one-to-two percent range;
  \item mean-shift $R^2$ is nonnegative or improves without a compensating
        degradation in Wasserstein distance;
  \item performance is stable across Brownian seeds and particle counts;
  \item learned diffusion is nonzero, bounded, and improves over both
        drift-only and pure-diffusion controls;
  \item improvement persists within major cell types rather than arising only
        from global mixture smoothing.
\end{itemize}

If population-level SDE training still matches persistence, the appropriate
conclusion is that expression snapshots and the present state representation do
not identify the requested dynamics at this resolution. Further architecture
sweeps would then be lower priority than lineage constraints, explicit cell
types, perturbations, or direct proliferation information.

\begin{thebibliography}{99}

\bibitem{lipman2023flow}
Lipman, Y., Chen, R. T. Q., Ben-Hamu, H., Nickel, M. \& Le, M.
Flow Matching for Generative Modeling.
\textit{International Conference on Learning Representations} (2023).
\url{https://arxiv.org/abs/2210.02747}

\bibitem{song2021sde}
Song, Y., Sohl-Dickstein, J., Kingma, D. P., Kumar, A., Ermon, S. \& Poole, B.
Score-Based Generative Modeling through Stochastic Differential Equations.
\textit{International Conference on Learning Representations} (2021).
\url{https://arxiv.org/abs/2011.13456}

\bibitem{li2020neuralsde}
Li, X., Wong, T.-K. L., Chen, R. T. Q. \& Duvenaud, D.
Scalable Gradients for Stochastic Differential Equations.
\textit{International Conference on Artificial Intelligence and Statistics} (2020).
\url{https://arxiv.org/abs/2001.01328}

\bibitem{debortoli2021diffusion}
De Bortoli, V., Thornton, J., Heng, J. \& Doucet, A.
Diffusion Schrodinger Bridge with Applications to Score-Based Generative Modeling.
\textit{Advances in Neural Information Processing Systems} (2021).
\url{https://arxiv.org/abs/2106.01357}

\bibitem{albergo2023stochastic}
Albergo, M. S., Boffi, N. M. \& Vanden-Eijnden, E.
Stochastic Interpolants: A Unifying Framework for Flows and Diffusions.
\url{https://arxiv.org/abs/2303.08797}

\bibitem{tong2024simulationfree}
Tong, A. et al.
Simulation-Free Schrodinger Bridges via Score and Flow Matching.
\textit{International Conference on Artificial Intelligence and Statistics} (2024).
\url{https://arxiv.org/abs/2307.03672}

\end{thebibliography}

\end{document}
